% Source: https://artofproblemsolving.com/community/c3h1728070_rsm_geometry_homework
% Stable identity: aops-topic-1728070-q01
Since $AB=AC$, the base angles are equal: $\angle ACB=\angle ABC=70^\circ$.

\begin{asy}
usepackage("amsmath");
size(5cm);
defaultpen(fontsize(10pt));
pair B=(0,0),C=(4,0),A=(2,2*Tan(70));
pair Ep=extension(C,C+dir(145),B,A);
pair D=extension(Ep,Ep+(1,0),A,C);
draw(A--B--C--cycle,linewidth(0.9));
draw(C--Ep,linewidth(0.8));draw(Ep--D,blue+linewidth(0.9));
draw(arc(C,0.8,145,180),linewidth(0.6));draw(arc(C,0.95,110,145),linewidth(0.6));
draw(arc(Ep,0.9,-35,0),linewidth(0.6));
label("$35^\circ$",C+1.3*dir(164),fontsize(8pt));
label("$35^\circ$",C+1.55*dir(128),fontsize(8pt));
label("$35^\circ$",Ep+1.5*dir(-17.5),fontsize(8pt));
draw(arc(B,0.6,0,70),linewidth(0.6));label("$70^\circ$",B+1.05*dir(35),fontsize(8pt));
path tk(pair P,pair Q){pair M=(P+Q)/2,u=unit(Q-P)*0.18,v=rotate(90)*u;return (M-v)--(M+v);}
draw(tk(C,D));draw(tk(Ep,D));
label("$A$",A,N);label("$B$",B,SW);label("$C$",C,SE);label("$E$",Ep,W);label("$D$",D,E);
label("\mbox{}",(min(currentpicture,user=true)+max(currentpicture,user=true))/2,black); // keeps the text after the figure black when LaTeX typesets the labels inline
\end{asy}

As $E$ lies on side $AB$, the ray $CE$ lies inside $\angle ACB$, so
\[
\angle DCE=\angle ACE=\angle ACB-\angle ECB=70^\circ-35^\circ=35^\circ .
\]
Triangle $CDE$ is isosceles with $CD=DE$, so its base angles at $C$ and $E$ are equal: $\angle DEC=\angle DCE=35^\circ$.

Hence $\angle DEC=\angle ECB=35^\circ$. These are alternate angles formed by the transversal $EC$ with the lines $DE$ and $BC$ ($D$ and $B$ lie on opposite sides of line $CE$). Equal alternate angles give $DE\parallel BC$. $\blacksquare$

(So $CE$ bisects $\angle ACB$, and the result is the general fact that the bisector from $C$ makes $CD=DE$ exactly when $DE\parallel BC$.)
